% Converted from README.md
\subsection{Primer semantičke greške}\label{primer-semantiux10dke-greux161ke}

Pogledajmo primer neispravne funkcije za proveru parnosti celog broja:

\begin{minted}{c}
int bad_even(int x) {
    if (x % 2 == 0)
        return 1;
    if (x % 2 == 1)
        return 0;
    return 1;
}
\end{minted}

Ovaj primer, za razliku od prethodnog ima grešku prilikom provere neparnih negativnih brojeva jer je za njih ostatak pri deljenju jednak \texttt{-1} pa nisu obuhvaćeni drugom \texttt{if} naredbom.

Pokušaćemo da korišćenjem KLEE alata otkrijemo ovu grešku i dobijemo test primer za koji je dostižna poslednja \texttt{return} naredba. Kao i malopre, u glavnom programu označavamo simboličku promenljivu i dodajemo poziv funkcije \texttt{klee\_assert(0)} koja će signalizirati grešku pri eksploatisanju putanje u kojoj je malopređašnja nedostižna naredba.

Nakon dobijanja LLVM bitkoda \texttt{bad\_even.bc} i pokretanja KLEE alata:

\begin{minted}{text}
KLEE: output directory is "klee-out-0"
Using STP solver backend
KLEE: ERROR: bad_even.c:12: ASSERTION FAIL: 0
KLEE: NOTE: now ignoring this error at this location

KLEE: done: total instructions = 39
KLEE: done: completed paths = 3
KLEE: done: generated tests = 3
\end{minted}

Primetimo da su sada pronađene tri putanje i generisani odgovarajući test primeri, ali je i prijavljena greška koju je signalizirala funkcija \texttt{klee\_assert}. Ukoliko pogledamo sadržaj direktorijuma sa test primerima, primetićemo da je napravljena datoteka \texttt{test000002.assert.err} koja nam ukazuje da je gore pomenuta greška signalizirana pri drugom test primeru.

Pogledajmo kakav je ulaz u pitanju:

\begin{minted}{bash}
$ ktest-tool --write-ints test000002.ktest
\end{minted}

\begin{minted}{text}
ktest file : 'test000002.ktest'
args       : ['bad_even.bc']
num objects: 1
object    0: name: 'x'
object    0: size: 4
object    0: data: -2147483647
\end{minted}

KLEE je pronašao da je sporna vrednost ulaza \texttt{-2147483647}. Upotrebom debagera i korišćenjem sporne ulazne vrednosti može se uočiti mesto problema, odnosno da je zaboravljeno da je ostatak pri deljenju sa \texttt{2} u slučaju negativnih brojeva \texttt{-1}, a ne \texttt{1}. Očekivano, vreme izvršavanja koje je potrošeno na solver je veće nego u slučaju ispravnog programa.

\begin{minted}{bash}
$ klee-stats klee-last
\end{minted}

\begin{minted}{text}
------------------------------------------------------------------------
|  Path   |  Instrs|  Time(s)|  ICov(%)|  BCov(%)|  ICount|  TSolver(%)|
------------------------------------------------------------------------
|klee-last|      39|     0.17|    11.30|     7.14|     292|       93.78|
------------------------------------------------------------------------
\end{minted}

Generalno, prikazanom upotrebom funkcije \texttt{klee\_assert} možemo označiti putanju za koju smo zainteresovani, bilo da bismo proverili da li je ta putanja izvodljiva ili da bismo imali na izlazu označen ulaz kojim se ta putanja dostiže.
